% GDZ PPN236018647,native559upper and its firstthreefootnotes;HKUST560.
\clearpage\SourceLeaf{559}\ArithmeticHead{530}{L}{TAGEBUCH MIT ERLÄUTERUNGEN.}
\NumberedItem{86}
\begingroup\renewcommand{\thefootnote}{[*)]}
\indent Theorema \textsc{la Grange} de transformatione functionum\footnotemark\ ad functiones quotcunque variabilium extendi.
\renewcommand{\thefootnote}{[*)}
\footnotetext{In der Handschrift steht \textls[120]{functionem}.]}
\par\noindent\hfill Gotting[ae, 1798] Mai.
\par\endgroup

\medskip\begingroup\small
\renewcommand{\thefootnote}{**)}
\indent Im Nachlaß findet sich keine Aufzeichnung, die hier in Betracht kommen könnte, aber in dem oben S. 429 abgedruckten Briefe an \textsc{Hindenburg} vom 8. Oktober 1799 erwähnt \textsc{Gauss} die hier angezeigte Verallgemeinerung der Umkehrungsformel von \textsc{Lagrange}\footnote{Daß es sich hier wirklich um den Umkehrungssatz von \textsc{Lagrange} handelt, wird dadurch außer Zweifel gesetzt, daß, wie die Handschrift erkennen läßt, \textsc{Gauss} ursprünglich statt \textls[120]{transformatione} schreiben wollte \textls[120]{reversione}.} und bemerkt, daß er inzwischen das Verfahren von \textsc{Laplace} zum Beweise dieser Formel, das ihm im April 1797 noch unbekannt war, kennen gelernt habe. Nun enthält der VIII. Abschnitt derselben Abhandlung von \textsc{Laplace}, in deren VII. Abschnitt jenes Verfahren entwickelt ist (siehe die zweite Fußnote auf S. 430), eine Verallgemeinerung des \textsc{Lagrange}schen Satzes auf Funktionen von zwei und mehr Veränderlichen\begingroup\renewcommand{\thefootnote}{***)}\footnote{Eine vereinfachte Darstellung dieser Verallgemeinerung von \textsc{Laplace} gibt G. \textsc{Darboux}, \textit{Sur la série de Laplace}, Comptes Rendus de l'Acad. des Sciences \textbf{68}, Paris 1869, S. 324, vergl. auch T.-J. \textsc{Stieltjes}, \textit{Sur une généralisation de la série de Lagrange}, Annales de l'École Normale (3) \textbf{2}, 1885, S. 93.}\endgroup. \textsc{Gauss} hat also bei seiner Verallgemeinerung jedenfalls an \textsc{Laplace} angeknüpft. Es ist dies darum besonders bemerkenswert, weil \textsc{Gauss} um diese Zeit, also etwa in der ersten Hälfte des Jahres 1798, überhaupt angefangen hat, sich mit den Schriften von \textsc{Laplace} zu beschäftigen, vergl. die Notiz Nr. 88 und die zugehörigen Bemerkungen.
\par\noindent\hfill\textsc{Schlesinger.}
\par\endgroup
